% =====================================================================
%  Theoretical model of the common-path Fresnel-biprism low-coherence
%  digital holography spectroscopy (FB-DHS) instrument
%  Proyecto Avanzado 1 -- A. F. Contreras (EAFIT)
%
%  Notation follows the course notes "Optica de Fourier" (Goodman-style):
%  complex amplitude U, object plane (xi, eta), pupil plane (x, y),
%  image plane (u, v), G = F{g}, h (x) U_g, t_l, P(x,y).
% =====================================================================
\documentclass[11pt,a4paper]{article}

\usepackage[utf8]{inputenc}
\usepackage[T1]{fontenc}
\usepackage{lmodern}
\usepackage[margin=2.5cm]{geometry}
\usepackage{amsmath,amssymb}
\usepackage{bm}
\usepackage{booktabs}
\usepackage{enumitem}
\usepackage[hidelinks]{hyperref}

% ---------- notation ----------
\newcommand{\vr}{\mathbf{r}}
\newcommand{\vq}{\mathbf{q}}
\newcommand{\shat}{\hat{\mathbf{s}}}
\newcommand{\FT}{\mathcal{F}}
\newcommand{\um}{\,\mu\mathrm{m}}
\newcommand{\nm}{\,\mathrm{nm}}
\newcommand{\mm}{\,\mathrm{mm}}
\newcommand{\ee}{e}
\newcommand{\dd}{d}
\newcommand{\avg}[1]{\left\langle #1 \right\rangle}
\newcommand{\rect}{\operatorname{rect}}
\newcommand{\sinc}{\operatorname{sinc}}
\newcommand{\step}{\operatorname{step}}
\newcommand{\comb}{\operatorname{comb}}

\title{Theoretical Model of the Common-Path Fresnel-Biprism\\
Digital Holography Spectroscopy (FB-DHS) Instrument}
\author{Andr\'es Felipe Contreras\\ \small Proyecto Avanzado 1 -- EAFIT University}
\date{Working draft -- \today}

\begin{document}
\maketitle

\begin{abstract}
This document derives the forward optical model that maps the scatterer
positions of a simulated sample to the fluctuation spectrum measured by a
common-path Fresnel-biprism digital holography instrument. The instrument
modeled is the EAFIT implementation built by L\'opez Zapata
\cite{lopez2026}, which follows the architecture of Goodwin et al.\
\cite{goodwin2025}. All numerical values refer to the EAFIT setup unless
stated otherwise. The model is scalar and paraxial, and it uses single
scattering. It is split into two
halves: \emph{illumination}, from the source to the field incident on each
scatterer, and \emph{collection}, from the scattered field to the camera,
the reconstruction and the spectrum. The objective sample is provolone
cheese, a passive emulsion of fat globules ($1$--$10\um$) in a casein matrix.
The notation follows the Fourier-optics course notes.
\end{abstract}

\tableofcontents

% =====================================================================
\section{Notation}
\label{sec:notation}

\subsection{Fourier optics conventions}

The conventions are those of the course notes, which follow Goodman \cite{goodman}.

\begin{itemize}[leftmargin=*]
  \item A monochromatic field of optical frequency $\nu$ is
        $u(P,t) = \mathrm{Re}\{U(P)\,\ee^{-i2\pi\nu t}\}$. The complex amplitude
        $U(P)$ satisfies the Helmholtz equation $(\nabla^2 + k^2)U = 0$, with
        $k = 2\pi/\lambda$, and the intensity is $I = |U|^2$.
  \item Two-dimensional Fourier transform and its inverse:
  \begin{align}
    G(f_x,f_y) &= \FT\{g(x,y)\} = \iint_{-\infty}^{\infty} g(x,y)\,
       \ee^{-i2\pi(f_xx + f_yy)}\,\dd x\,\dd y,\\
    g(x,y) &= \FT^{-1}\{G\} = \iint_{-\infty}^{\infty} G(f_x,f_y)\,\ee^{\,i2\pi(f_xx+f_yy)}\,\dd f_x\,\dd f_y .
  \end{align}
  The theorems used below are the shift theorem
  $\FT\{g(x-a,y-b)\} = G\,\ee^{-i2\pi(f_xa+f_yb)}$, the convolution theorem
  $\FT\{g\otimes h\} = GH$, and the correlation theorem
  $\FT\{\iint g(\xi,\eta)\,h^*(\xi - x,\eta - y)\,\dd\xi\,\dd\eta\} = G^*H$
  (it reduces to $|G|^2$ for the autocorrelation).
  \item Plane waves are written with direction cosines $(\alpha,\beta,\gamma)$,
        where $\gamma = \sqrt{1-\alpha^2-\beta^2}$, as
        $U = A\,\ee^{\,ik(\alpha x + \beta y + \gamma z)}$.
  \item Fresnel diffraction over a distance $z$:
  \begin{equation}
    U(x,y,z) = \frac{\ee^{ikz}}{i\lambda z}\iint U(\xi,\eta,0)\,
       \ee^{\,i\frac{k}{2z}\left[(x-\xi)^2+(y-\eta)^2\right]}\,\dd\xi\,\dd\eta,
    \qquad
    H(f_x,f_y) = \ee^{ikz}\,\ee^{-i\pi\lambda z(f_x^2+f_y^2)}.
    \label{eq:fresnel}
  \end{equation}
  \item Thin element of thickness function $\Delta(x,y)$ and index $n$:
        $t(x,y) = \ee^{\,ik\Delta_0}\,\ee^{\,ik(n-1)\Delta(x,y)}$. For a thin lens,
        $t_l(x,y) = \ee^{-i\frac{k}{2f}(x^2+y^2)}$ (constant phase dropped),
        and the pupil function $P(x,y)$ is 1 inside the aperture and 0 outside.
  \item Coherent imaging with magnification $M$:
  \begin{gather}
    U_i(u,v) = h(u,v)\otimes U_g(u,v),
    \qquad
    U_g(\tilde\xi,\tilde\eta) = \frac{1}{|M|}\,U_o\!\left(\frac{\tilde\xi}{M},\frac{\tilde\eta}{M}\right),
    \label{eq:imaging}\\
    h(u,v) = \frac{A}{\lambda^2 z_1 z_2}\iint P(x,y)\,
       \ee^{-i\frac{2\pi}{\lambda z_2}(ux+vy)}\,\dd x\,\dd y,
    \nonumber
  \end{gather}
  where $\tilde\xi = M\xi$, $\tilde\eta = M\eta$ and $U_g$ is the ideal
  (geometrical-optics) image.
\end{itemize}

\subsection{Planes and coordinates}

\begin{table}[ht]
\centering
\small
\caption{Planes of the instrument and their coordinates.}
\label{tab:planes}
\begin{tabular}{@{}lll@{}}
\toprule
Plane & Coordinates & Field \\
\midrule
Sample (surface at $z=0$, depth $z>0$) & $(\xi,\eta)$, $\vr = (\xi,\eta,z)$ & $U_o(\xi,\eta)$ \\
Objective pupil (L1) & $(x,y)$ & $P_1(x,y)$ \\
Intermediate image & $(u,v)$ & $U_i(u,v)$ \\
Fourier lens (L2) & $(x',y')$ & $P_2(x',y')$ \\
Biprism, a distance $d$ before the camera & $(x_B,y_B)$ & $U_B$ \\
Camera = back focal plane of L2 & $(x_c,y_c)$ & $U_c(x_c,y_c)$, $I_H(x_c,y_c)$ \\
Reconstruction (lag in image coordinates) & $(\Delta u,\Delta v)$ & $S_+(\Delta u,\Delta v)$ \\
\bottomrule
\end{tabular}
\end{table}

Two other symbols are needed. The temporal frequency of the intensity
fluctuations is written $f_t$, to keep it apart from the spatial frequencies
$(f_x,f_y)$ and from focal lengths. The light spectrum is written as a
function of the vacuum wavenumber $k$.

% =====================================================================
\section{Scope and assumptions}
\label{sec:scope}

\begin{enumerate}[label=(A\arabic*)]
  \item \textbf{Scalar, paraxial optics.} Polarization is ignored. Free-space
        propagation between optical elements uses the Fresnel approximation,
        Eq.~\eqref{eq:fresnel}.
  \item \textbf{Single scattering.} Each detected photon is scattered once by
        a globule. Between the surface and that globule, the light travels
        as the attenuated ballistic beam (first-order Born approximation
        with extinction). Multiple scattering is left for later work.
  \item \textbf{Quasi-static sample.} The globules move negligibly during one
        optical coherence time ($\sim 10^{-13}\,$s). So, for each wavenumber,
        the field is computed from the instantaneous positions $\vr_j(t)$.
  \item \textbf{Low temporal coherence, full spatial coherence.} The
        single-mode fiber gives a spatially coherent beam. Different
        wavenumbers $k$ are mutually incoherent, so the camera adds
        intensities over the spectrum.
  \item \textbf{Static scattering amplitudes.} Globule radii do not change,
        so each scattering amplitude is a constant. Only the phases depend on
        time.
\end{enumerate}

The input to the model is the set of globule positions $\{\vr_j(t_n)\}$ and
radii $\{a_j\}$ produced by the LAMMPS simulation at the camera frame times.
The output is the normalized fluctuation power spectrum $\Phi(f_t)$.

% =====================================================================
\section{The EAFIT instrument}
\label{sec:params}

The model describes the common-path setup of L\'opez Zapata
\cite[\S3.1.4--3.1.6, Fig.~3.9]{lopez2026}. The sample sits in a multi-well
plate and is illuminated from below by the SLD (inverted configuration). The
backscattered light is folded by a mirror into a horizontal collection arm.
There, an objective lens (OL) in a $2f$--$2f$ configuration forms a real image
on the image plane (IP), and a Fourier lens (FL) placed one focal length after
the IP puts the camera on its back focal plane (FP). The Fresnel biprism (FB)
sits between FL and the camera. An adjustable aperture before FL sets the beam
diameter $D$ at the biprism, and Fig.~3.9 of Ref.~\cite{lopez2026} also shows
an aperture at the IP, which acts as a field stop.

The same architecture was introduced by Goodwin et al.\ \cite{goodwin2025},
whose provolone measurements are the main external reference. Their values are
listed for comparison only. Every quantity in Table~\ref{tab:params} stays a
parameter of the model.

\begin{table}[ht]
\centering
\small
\setlength{\tabcolsep}{4pt}
\caption{Instrument parameters. The EAFIT column is the one used in this
model. Entries marked $^\dagger$ are inferred here, not stated in the source.}
\label{tab:params}
\begin{tabular}{@{}lll@{}}
\toprule
Parameter & EAFIT \cite{lopez2026} (this model) & Goodwin et al.\ \cite{goodwin2025} (reference) \\
\midrule
Source & Thorlabs SLD850S-A10W & Superlum SLD-381-HP2 \\
Center wavelength $\lambda_0$ & $850\nm$ & $839.6\nm$ \\
Bandwidth (FWHM) $\Delta\lambda$ & $60\nm$ & $27.7\nm$ \\
Power & $\approx 9.3\,$mW (source) & $16.5\,$mW (at sample) \\
Sample holder & multi-well plate, lit from below & 96-well plate, lit from below \\
Illumination angle $\psi$ (in air) & oblique in Fig.~3.9, value not stated & $30^\circ$ \\
Objective (OL) & $f_1 = 75\mm$, $2f$--$2f$, $M=-1$ & not given \\
Fourier lens (FL) & $f_2 = 200\mm$ & not given \\
Biprism & apex $170^\circ$ ($\beta = 5^\circ$) & not given \\
Camera & DMK 41BU02, $p = 3.75\um$ & FLIR BFS-U3-32S4M, $p=3.45\um$ \\
Hologram ROI & $800\times800$ px (Spirulina data) & $800\times800$ px \\
Frames, rate & 1400 at $f_s = 0.2\,$Hz ($T\approx7000\,$s) & 3000 at 3 fps \\
Spectral range $[1/T, f_s/2]$ & $0.14\,$mHz -- $0.1\,$Hz & $0.33\,$mHz -- $1.5\,$Hz \\
Exposure $T_\mathrm{exp}$ & not stated for provolone & $\le 300\,$ms \\
Coherence length $l_c$, Eq.~\eqref{eq:lc} & $5.3\um$ & $11.2\um$ \\
Fringe period $\Lambda$, Eq.~\eqref{eq:biprism} & $9.6\um \approx 2.6\,p^\dagger$ & $\approx 3p$ \\
Field of view (rule of nine) & $\approx 5\mm^\dagger$ & $\approx 2\mm$ (beam profile) \\
Spectrum observable & $|\FT_t\{S_+\}|^2$ (field) & $|S_+|^2$ fluctuations$^\dagger$ \\
\bottomrule
\end{tabular}
\end{table}

\paragraph{Sample.}
Fat volume fraction $\phi \approx 0.3$ (USDA composition). Refractive indices
at $\approx 840\nm$: fat $n_p \approx 1.455$, casein matrix
$n_m \approx 1.392$ \cite{stocker2016}, so the relative index is
$m = n_p/n_m \approx 1.045$. Globule radii range over $a \approx 0.5$--$5\um$.

% =====================================================================
\section{Part I: Illumination (source to scatterer)}
\label{sec:illum}

\subsection{Source spectrum and temporal coherence}

The source field is a superposition of mutually incoherent monochromatic
components $u(P,t) = \mathrm{Re}\{U(P;k)\,\ee^{-i2\pi\nu t}\}$, each with vacuum
wavenumber $k = 2\pi\nu/c$ and weight $S(k)$. The SLD spectrum is modeled as a
Gaussian,
\begin{equation}
  S(k) = \frac{1}{\sqrt{2\pi}\,\sigma_k}
         \exp\!\left[-\frac{(k-k_0)^2}{2\sigma_k^2}\right],
  \qquad
  \sigma_k = \frac{\Delta k}{2\sqrt{2\ln 2}},
  \qquad
  \Delta k = \frac{2\pi\,\Delta\lambda}{\lambda_0^2},
  \label{eq:spectrum}
\end{equation}
normalized so that $\int S\,\dd k = 1$. Two fields that reach the camera with
an optical path difference $\Delta L$ interfere with the weight given by the
complex degree of temporal coherence, which is the Fourier transform of the
spectrum (Wiener--Khinchin):
\begin{equation}
  \gamma(\Delta L) = \int S(k)\,\ee^{\,ik\Delta L}\,\dd k
  = \exp\!\left(-\frac{\sigma_k^2\,\Delta L^2}{2}\right)\ee^{\,ik_0\Delta L}.
  \label{eq:gamma}
\end{equation}
The envelope $|\gamma|$ falls to one half at $|\Delta L| = l_c$, where
\begin{equation}
  l_c = \frac{\sqrt{2\ln 2}}{\sigma_k} = \frac{2\ln 2}{\pi}\,\frac{\lambda_0^2}{\Delta\lambda}.
  \label{eq:lc}
\end{equation}
For the EAFIT source this gives $\sigma_k \approx 0.22\um^{-1}$ and
$l_c \approx 5.3\um$ (the Goodwin source gives $11.2\um$). Ref.~\cite{lopez2026}
quotes $\lambda_0^2/\Delta\lambda \approx 12\um$, which is a different
convention for the same quantity.

\paragraph{Spectral discretization.}
In the simulation, $S(k)$ is sampled at $N_k$ points with spacing $\delta k$
over $k_0 \pm 3\sigma_k$. The sampled coherence function repeats with period
$2\pi/\delta k$ in $\Delta L$, so this period must exceed the largest path
difference in the sample, $\Delta L_\mathrm{max} \approx 2n_m z_\mathrm{max}$:
\begin{equation}
  \delta k < \frac{2\pi}{2\,\Delta L_\mathrm{max}}, \qquad
  N_k \approx \frac{6\sigma_k}{\delta k}.
  \label{eq:Nk}
\end{equation}
For a slab $z_\mathrm{max} = 20\um$, $\Delta L_\mathrm{max}\approx 56\um$,
which gives $N_k \approx 24$ for the EAFIT source.

\subsection{Beam at the sample}

The SLD output reaches the sample as a Gaussian beam with waist $w$. The
beam size is not stated in Ref.~\cite{lopez2026}. The field of view is set by
whichever is smaller: the beam, or the field stop at the IP (about $5\mm$
under the rule of nine, Section~\ref{sec:collect}). For $w \sim 1\mm$ the
Rayleigh range $z_R = \pi w^2/\lambda$ is of order metres, much larger than the
sample depth. So, inside the sample,
the illumination is a plane wave with a Gaussian transverse envelope. With
direction cosines $(\alpha_t,\beta_t,\gamma_t)$ in the medium,
\begin{equation}
  U_\mathrm{inc}^{(0)}(\xi,\eta,z;k) = \sqrt{S(k)}\,A\,
     \exp\!\left(-\frac{\xi^2+\eta^2}{w^2}\right)
     \exp\!\left[\,ik\,n_m(\alpha_t\xi + \beta_t\eta + \gamma_t z)\right].
  \label{eq:beam}
\end{equation}

\subsection{Geometry and refraction}

In the EAFIT setup the collection axis is the sample normal: the fold mirror
under the plate sends the vertically backscattered light into the horizontal
collection arm. Fig.~3.9 of Ref.~\cite{lopez2026} draws the illumination
reaching the plate obliquely from below, but the angle is not stated, so it is
kept as a parameter $\psi$ (in air, in the $\xi z$ plane). The plate bottom is
a parallel slab, so it does not change the angle inside the sample, and
Snell's law gives the refracted angle $\theta_t$:
\begin{equation}
  \sin\theta_t = \frac{\sin\psi}{n_m},
  \qquad
  (\alpha_t,\beta_t,\gamma_t) = (\sin\theta_t,\,0,\,\cos\theta_t).
  \label{eq:snell}
\end{equation}
For $\psi = 30^\circ$ (Goodwin's value), $\theta_t \approx 21.1^\circ$. For
on-axis illumination, $\theta_t = 0$. The Fresnel transmission losses at the
plate bottom only rescale $A$ and are ignored.

The light collected by the objective leaves the sample with direction cosines
close to $(0,0,-1)$. Writing $\shat_t = (\alpha_t,\beta_t,\gamma_t)$ and
$\shat_\mathrm{out} = (0,0,-1)$, the scattering vector is
\begin{equation}
  \vq = k\,n_m\,(\shat_t - \shat_\mathrm{out}) = k\,n_m\,(\sin\theta_t,\;0,\;1+\cos\theta_t),
  \qquad
  |\vq| = 2k\,n_m\cos(\theta_t/2).
  \label{eq:q}
\end{equation}
At $\lambda_0 = 850\nm$, $|\vq| = 4\pi n_m/\lambda_0 \approx 20.6\um^{-1}$ for
$\psi = 0$ and $|\vq| \approx 20.2\um^{-1}$ for $\psi = 30^\circ$. The
difference is below $2\%$, so $q \approx 20.5\um^{-1}$ is used for estimates.
The angle matters more for the direction of $\vq$, which sets the tilt of the
coherence gate (Section~\ref{sec:collect}). The in-medium scattering angle is
$\theta_s = \pi - \theta_t$.

\subsection{Ballistic attenuation}

Under (A2), the beam reaching a globule at depth $z$ has travelled a path
$z/\cos\theta_t$ through the turbid medium, and the scattered field returns
along a path $z$. The ballistic amplitude decays as
\begin{equation}
  T(z) = \exp\!\left[-\frac{\mu_t z}{2}\left(\frac{1}{\cos\theta_t} + 1\right)\right],
  \qquad \mu_t = \mu_s + \mu_a \approx \mu_s,
  \label{eq:atten}
\end{equation}
since absorption is negligible in the near infrared. The scattering
coefficient follows from the number density
$n_v = \phi/(\tfrac{4}{3}\pi a^3)$ and the scattering efficiency $Q_s$:
\begin{equation}
  \mu_s = n_v\,\pi a^2 Q_s = \frac{3\phi\,Q_s}{4a}.
\end{equation}
For $m - 1 \ll 1$ and $ka \gg 1$, the anomalous-diffraction approximation
gives
\begin{equation}
  Q_s(\varrho) = 2 - \frac{4}{\varrho}\sin\varrho + \frac{4}{\varrho^2}(1-\cos\varrho),
  \qquad \varrho = 2x_s(m-1), \quad x_s = k\,n_m a,
\end{equation}
where $x_s$ is the size parameter. For $a = 2\um$, $x_s \approx 21$,
$\varrho \approx 1.9$, $Q_s \approx 1.45$, and $\mu_s \approx 0.16\um^{-1}$, so
the scattering mean free path is $\ell_s \approx 6\um$. At $\phi \approx 0.3$
the globules are densely packed, and a hard-sphere structure factor
(Percus--Yevick) would lower this value. It should be treated as an upper
bound.

\paragraph{Consequence for the sample size.}
Single scattering is only self-consistent within the top one or two $\ell_s$.
A slab of depth $z_\mathrm{max} \approx 15$--$20\um$ under a free surface
contains essentially all of the signal ($T^2 < 10^{-2}$ beyond that depth).

\subsection{Result of Part I: field incident on globule \texorpdfstring{$j$}{j}}

Combining Eqs.~\eqref{eq:beam} and \eqref{eq:atten}, the field that
illuminates globule $j$, at $\vr_j = (\xi_j,\eta_j,z_j)$, is
\begin{equation}
  \boxed{\;
  U_\mathrm{inc}(\vr_j;k,t) = \sqrt{S(k)}\;A\;
     \ee^{-(\xi_j^2+\eta_j^2)/w^2}\;
     \ee^{-\mu_t z_j/(2\cos\theta_t)}\;
     \ee^{\,ik\,n_m\,\shat_t\cdot\vr_j(t)}
  \;}
  \label{eq:Einc}
\end{equation}
The last factor is the illumination half of the Doppler phase. The
collection half, $-k\,n_m\,\shat_\mathrm{out}\cdot\vr_j = k\,n_m z_j$, is added in
Part~II, and together the two halves make $\vq\cdot\vr_j$.

% =====================================================================
\section{Scattering by a single globule}
\label{sec:scatter}

Globule $j$ re-radiates the incident field with the far-field amplitude
$S_1(\theta_s; a_j, m, x_s)$ from Mie theory (scalar version: the average of
the two polarization amplitudes). The size parameter is
$x_s \approx 10\,a[\mu\mathrm{m}]$, so the globules are in the Mie regime. The
Rayleigh--Gans--Debye approximation needs $\varrho \ll 1$, which only holds
for $a \lesssim 0.5\um$. In practice, a table $a_s(a) \equiv S_1(\theta_s; a)$
is computed once over the radius distribution at the instrument's $\theta_s$.

Seen from the objective, which images the plane $z=0$, each globule acts as a
point source on the sample plane, with complex strength
\begin{equation}
  c_j(k,t) = a_s(a_j)\;T(z_j)\;\ee^{-(\xi_j^2+\eta_j^2)/w^2}\;
             \ee^{\,i\vq\cdot\vr_j(t)} .
  \label{eq:cj}
\end{equation}
Common phase factors that cancel in every interference term have been
dropped. The object field is then
\begin{equation}
  U_o(\xi,\eta;k,t) = \sum_j c_j(k,t)\,\delta(\xi-\xi_j,\,\eta-\eta_j).
  \label{eq:Uo}
\end{equation}
Because $|a_s|^2$ grows steeply with radius, the largest globules dominate the
signal. The defocus of globules below $z=0$ is neglected; this holds while
$z_\mathrm{max}$ is smaller than the depth of field $\lambda/\mathrm{NA}^2$.

% =====================================================================
\section{Part II: Collection (scatterer to camera)}
\label{sec:collect}

\subsection{Coherent imaging by the objective}

Inserting Eq.~\eqref{eq:Uo} into the imaging relation Eq.~\eqref{eq:imaging},
and using $\delta(\tilde\xi/M - \xi_j) = |M|\,\delta(\tilde\xi - M\xi_j)$, the
ideal image is a set of point sources at $(M\xi_j, M\eta_j)$, and the image
field is a sum of shifted impulse responses:
\begin{gather}
  U_i(u,v;k,t) = |M|\sum_j c_j(k,t)\;h_1(u - M\xi_j,\; v - M\eta_j),
  \label{eq:Ui}\\
  h_1(u,v) = \frac{A}{\lambda^2 z_1 z_2}\iint P_1(x,y)\,
       \ee^{-i\frac{2\pi}{\lambda z_2}(ux+vy)}\,\dd x\,\dd y .
  \nonumber
\end{gather}
The width of $h_1$ referred to the sample is the resolution cell,
$\delta \approx \lambda/(2\,\mathrm{NA})$. This is the same pupil model as
Ref.~\cite[Eqs.~2.7--2.10]{lopez2026}, where the coherent transfer function is
the pupil with cutoff $f_c$. For the EAFIT objective ($f_1 = 75\mm$ in
$2f$--$2f$), $z_1 = z_2 = 150\mm$ and $M = -1$. The lens diameter is not
stated. For a 1-inch lens, $\mathrm{NA} \approx 12.7/150 \approx 0.085$ and
$\delta \approx 5\um$ (inferred).

\paragraph{Field stop at the image plane.}
The aperture drawn at the IP in Fig.~3.9 of Ref.~\cite{lopez2026} multiplies
the image by its pupil function, $U_i \to U_i\,P_\mathrm{IP}(u,v)$. Its diameter
$D_\mathrm{IP}$ sets the field of view at the sample, $L = D_\mathrm{IP}/|M|$
(unless the illumination beam is smaller), and therefore the speckle size on
the camera.

\subsection{Fourier lens}

The intermediate image sits at the front focal plane of L2. Using the
lens-as-Fourier-transform result with $d = f_2$ (and ignoring vignetting by
$P_2$), the field on the back focal plane is an exact Fourier transform:
\begin{equation}
  U_F(x_c,y_c;k,t) = \frac{1}{i\lambda f_2}\,G_i\!\left(\frac{x_c}{\lambda f_2},\frac{y_c}{\lambda f_2}\right),
  \qquad
  G_i(f_u,f_v) = \FT\{U_i(u,v)\}.
  \label{eq:UF}
\end{equation}
With Eq.~\eqref{eq:Ui} and the shift and convolution theorems, the Fourier
field is the scaled pupil, filled with a random sum of plane-wave phases:
\begin{equation}
  G_i(f_u,f_v) = |M|\,H_1(f_u,f_v)\sum_j c_j\,
     \ee^{-i2\pi M(f_u\xi_j + f_v\eta_j)},
  \label{eq:Gi}
\end{equation}
with $H_1 = \FT\{h_1\} \propto P_1(-\lambda z_2 f_u,\,-\lambda z_2 f_v)$.
Two length scales appear on the camera. The first is the diameter $D_F$ of the
illuminated disc. The objective alone would give $2f_2\,\mathrm{NA}/|M|
\approx 34\mm$, but the adjustable aperture before FL limits the beam to a
smaller $D$, so $D_F = D$ in practice. The second is the speckle grain inside
that disc,
\begin{equation}
  s \approx \frac{\lambda f_2}{D_\mathrm{IP}} = \frac{\lambda f_2}{|M|L},
  \label{eq:speckle}
\end{equation}
set by the field of view, because the phases in the sum vary on the scale
$\Delta f_u \sim 1/(|M|L)$.

\subsection{Fresnel biprism: shear and carrier}

The biprism has apex angle $\alpha_B$, base angle $\beta = (\pi-\alpha_B)/2$
and index $n_g$. Its apex line is perpendicular to the unit vector
$\hat{\mathbf{e}}$, which is rotated by $\approx 45^\circ$ with respect to the
pixel grid. Let $x_B$ be the coordinate along $\hat{\mathbf{e}}$. The thickness
is largest at the apex, $\Delta(x_B) = \Delta_0 - \beta|x_B|$, so the thin-element
transmittance is
\begin{align}
  t_B(x_B,y_B) &= \ee^{\,ik n_g\Delta_0}\,\ee^{-ik(n_g-1)\beta\,|x_B|}
  \nonumber\\
  &= \ee^{\,ikn_g\Delta_0}\left[\step(x_B)\,\ee^{-ik\theta x_B}
    + \step(-x_B)\,\ee^{\,ik\theta x_B}\right],
  \qquad \theta = (n_g-1)\beta .
  \label{eq:tB}
\end{align}
Each half acts as a prism that tilts its part of the beam by $\theta$ toward
the axis; in direction cosines, $\alpha = \mp\theta$ for $x_B \gtrless 0$.

\paragraph{Propagation to the camera.}
The biprism sits a distance $d$ before the back focal plane of L2. Let
$U_B(x_B,y_B)$ be the field arriving at it, so that propagating $U_B$ by $d$
without the biprism would give $U_F$. In terms of the Fresnel transfer
function, a linear phase $\ee^{\mp ik\theta x_B}$ shifts the angular spectrum by
$\mp\theta/\lambda$, and after a distance $d$ this becomes a lateral shift
$\mp\theta d$ plus the same linear phase (shift theorem applied to
Eq.~\eqref{eq:fresnel}). Writing $U_{F\pm}$ for the field that each half of the
biprism aperture alone would produce at the camera,
\begin{gather}
  U_c(x_c,y_c) = U_{F+}(x_c+\eta_s,\,y_c)\,\ee^{-iK x_c/2}
               + U_{F-}(x_c-\eta_s,\,y_c)\,\ee^{\,iK x_c/2},
  \label{eq:biprism}\\
  K = 2k\theta = \frac{2\pi}{\Lambda},
  \qquad \Lambda = \frac{\lambda}{2\theta},
  \qquad \eta_s = \theta d,
  \nonumber
\end{gather}
where $x_c$ is now measured along $\hat{\mathbf{e}}$ and a constant phase
$\ee^{-i\pi\lambda d(\theta/\lambda)^2}$ has been dropped. This is Eq.~(4) of
Ref.~\cite{goodwin2025}, written here as $\eta_s$ for the shear to avoid
confusion with the sample coordinate $\eta$. The sign of $K$ is opposite to
Ref.~\cite{goodwin2025} because of the $\ee^{-i2\pi\nu t}$ convention. This only
changes which sideband is labelled $+1$.

\paragraph{Half-plane approximation.}
If the biprism is close enough to L2 that the beam from each image point
still covers the full aperture, $U_{F\pm}$ is simply $U_F$ restricted to one
half of the disc:
\begin{equation}
  U_{F\pm}(x_c,y_c) \approx \step(\pm x_c)\,U_F(x_c,y_c).
\end{equation}
This assumes that the image field of view is small compared with the
aperture at the biprism. Diffraction at the apex line is neglected.

\paragraph{Overlap geometry and the EAFIT numbers.}
With $\eta_s = D/4$, the two half-discs land on top of each other over a
width $W \approx D/2$, which is the diamond-shaped region of
Refs.~\cite{goodwin2025,lopez2026} (their Eqs.~(2) and (2.23)--(2.24)):
$W \approx D/2$ at the distance $d \approx D/(4\theta)$. The design rule is that
$W$ matches the sensor size $C$ and $d$ matches the distance to the camera.
For the EAFIT values (inferred where marked):
\begin{itemize}[leftmargin=*]
  \item Biprism: $\beta = 5^\circ$ and $n_g \approx 1.51$ (BK7 at $850\nm$$^\dagger$)
        give $\theta \approx 44.5\,$mrad and $\Lambda = \lambda/(2\theta)
        \approx 9.6\um \approx 2.6\,p$, slightly under the rule-of-nine value
        of $3p$ but still above the Nyquist limit of $2p$.
  \item Field of view: the rule of nine, $s = 9p$ in Eq.~\eqref{eq:speckle},
        gives $D_\mathrm{IP} = \lambda f_2/(9p) \approx 5\mm$, so
        $L \approx 5\mm$ at the sample.
  \item Beam at the biprism: the DMK 41BU02 sensor is $1280\times960$ px
        $\approx 4.8\times3.6\mm$$^\dagger$, so $W \approx C$ asks for
        $D \approx 7$--$10\mm$ and $d \approx 40$--$55\mm$$^\dagger$.
  \item Reconstruction: with an $N = 800$ ROI, the reconstruction pixel is
        $\lambda f_2/(Np) \approx 57\um$, the sideband centre is
        $Np/\Lambda \approx 314$ px from the zero order, and the sideband
        (an autocorrelation, so twice the field of view wide) spans about
        $\pm 88$ px around it.
\end{itemize}

\paragraph{Shear as a scattering-vector offset.}
$U_{F+}(x_c+\eta_s)$ and $U_{F-}(x_c-\eta_s)$ sample $G_i$ at spatial
frequencies that differ by $2\eta_s/(\lambda f_2)$. Referred to the sample,
the two interfering fields come from slightly different scattering vectors,
$\vq_\pm = \vq \pm \tfrac{1}{2}\Delta\vq_\perp$, with
$|\Delta\vq_\perp| = k\,|M|\,2\eta_s/f_2 \approx k\,\mathrm{NA}$.

\subsection{Low-coherence hologram}

For each wavenumber, the camera records $|U_c|^2$. Under (A4), the hologram
is the spectral average
\begin{align}
  I_H(x_c,y_c,t) &= \int S(k)\,|U_c(x_c,y_c;k,t)|^2\,\dd k
  \nonumber\\
  &= \underbrace{\avg{|U_{F+}|^2 + |U_{F-}|^2}_k}_{\text{zero order}}
  + \underbrace{\avg{U_{F+}(x_c+\eta_s)\,U_{F-}^*(x_c-\eta_s)}_k\,\ee^{-iKx_c}}_{+1\ \text{sideband}}
  + \text{c.c.},
  \label{eq:hologram}
\end{align}
where $\avg{\cdot}_k$ denotes the $S(k)$-weighted integral. Inserting
Eqs.~\eqref{eq:cj} and \eqref{eq:Gi} into the cross term, each globule pair
$(j,j')$ carries the factor $\ee^{\,i\vq\cdot(\vr_j-\vr_{j'})}$, in which $\vq$ is
proportional to $k$. The $k$ integral turns it into $\gamma(\Delta L_{jj'})$,
with the single-scattering optical path difference
\begin{equation}
  \Delta L_{jj'} = \frac{1}{k}\,\vq\cdot(\vr_j - \vr_{j'}).
\end{equation}
So a pair contributes to the sideband only if
\begin{equation}
  \left|\hat{\vq}\cdot(\vr_j - \vr_{j'})\right| \lesssim \frac{l_c}{2n_m\cos(\theta_t/2)}.
  \label{eq:gate}
\end{equation}
This is the ``stacked'' coherence gate of Ref.~\cite{goodwin2025}. For the
EAFIT source it is a slab only about $1.9\um$ thick along $\hat{\vq}$
(about $4\um$ for Goodwin's source), tilted by $\theta_t/2$ from the surface. The reconstruction adds all such
slabs incoherently, weighted by $T^2(z)$.

\paragraph{Chromatic scaling.}
For a fixed camera position, the spatial frequency $x_c/(\lambda f_2)$
depends on $\lambda$, so speckle at the edge of the disc blurs radially over
the bandwidth. The carrier also scales, since $\Lambda \propto \lambda$.
However, the sideband position in the reconstruction, $\lambda f_2/\Lambda =
2\theta f_2$, does not depend on $\lambda$, and neither does the
reconstruction coordinate. A per-$k$ simulation therefore has to use the
correct $\lambda$ in Eqs.~\eqref{eq:UF} and \eqref{eq:biprism}, but all
wavenumbers reconstruct onto the same grid.

\subsection{Camera sampling and exposure}

The camera samples $I_H$ with pixel pitch $p$ on an $N\times N$ grid. Frame
$n$ integrates over each pixel and over the exposure:
\begin{equation}
  I_n(x_c,y_c) = \left[\bar I_{H,n} \otimes \rect\!\left(\frac{x_c}{p}\right)\rect\!\left(\frac{y_c}{p}\right)\right]
     \frac{1}{p^2}\comb\!\left(\frac{x_c}{p}\right)\comb\!\left(\frac{y_c}{p}\right),
  \label{eq:camera}
\end{equation}
where $\bar I_{H,n}(x_c,y_c) = \int_{t_n}^{t_n+T_\mathrm{exp}} I_H(x_c,y_c,t)\,\dd t$
and $t_n = n/f_s$. By the convolution theorem, the pixel integration
multiplies the reconstruction by $\sinc(p f_x)\sinc(p f_y)$, and the comb
replicates it with period $1/p$ (the sampling theorem). In time, the exposure
multiplies the fluctuation spectrum by $\sinc^2(f_tT_\mathrm{exp})$, and,
together with aliasing above the Nyquist frequency $f_s/2$, produces the
``strobe'' floor of Ref.~\cite{goodwin2025}. Shot noise and quantization
are then added. With the EAFIT protocol ($f_s = 0.2\,$Hz), the exposure is much
shorter than the 5 s frame period, so every dynamics faster than $0.1\,$Hz
aliases into the measured band instead of being averaged out. In the
simulation, the exposure integral needs LAMMPS positions at several sub-steps
within each exposure, and the run must cover the full $T \approx 7000\,$s.

\subsection{Reconstruction}

Because the camera sits on the Fourier plane, one 2D Fourier transform of
each frame returns to the image domain. Taking the transform of
Eq.~\eqref{eq:hologram} with respect to $(x_c,y_c)$ and evaluating it at
$(f_x,f_y) = (\Delta u,\Delta v)/(\lambda f_2)$, the zero order lands at the
center and the two sidebands land at
$\Delta u = \pm\lambda f_2/\Lambda = \pm 2\theta f_2$ along $\hat{\mathbf{e}}$.
A window $W_+$ isolates one sideband:
\begin{equation}
  S_+(\Delta u,\Delta v; t_n) = W_+(\Delta u,\Delta v)\cdot\FT\{I_n(x_c,y_c)\}\Big|_{f_x=\Delta u/\lambda f_2,\,f_y=\Delta v/\lambda f_2}.
\end{equation}
By the correlation theorem, the transform of the product
$U_{F+}U_{F-}^*$ is the cross-correlation of the two half-aperture images
$U_{i\pm} = \FT^{-1}\{U_{F\pm}\}$:
\begin{equation}
  S_+(\Delta u,\Delta v,t) \propto \iint U_{i+}(u,v,t)\,U_{i-}^*(u-\Delta u,\,v-\Delta v,\,t)\,\dd u\,\dd v .
  \label{eq:Scorr}
\end{equation}
This is the autocorrelation of the image of Ref.~\cite{goodwin2025}, Eq.~(6).
Each $U_{i\pm}$ is the image formed through one half of the pupil, so it has
the form of Eq.~\eqref{eq:Ui} with $h_1$ replaced by the half-pupil impulse
response $h_\pm$ (pupil $P_1\,\step(\pm x)$), and with the shear phase of
$\vq_\pm$. Substituting gives the sideband in terms of globules:
\begin{equation}
  \boxed{
  \begin{aligned}
  S_+(\Delta u,\Delta v,t) \approx \sum_{j,j'} &a_s(a_j)\,a_s^*(a_{j'})\,T_jT_{j'}\;
     C\!\left(\Delta u - M(\xi_j - \xi_{j'}),\;\Delta v - M(\eta_j-\eta_{j'})\right)\\
     &\times\left|\gamma(\Delta L_{jj'})\right|\;
     \ee^{\,i\left[\vq_+\cdot\vr_j(t) - \vq_-\cdot\vr_{j'}(t)\right]}
  \end{aligned}}
  \label{eq:Splus}
\end{equation}
where $C(u,v) = \iint h_+(u',v')\,h_-^*(u'-u,v'-v)\,\dd u'\,\dd v'$ is the
cross-correlation of the two half-pupil impulse responses. Each
reconstruction pixel is therefore a sum over globule pairs separated laterally
by $\approx(\Delta u,\Delta v)/M$ and lying in the same coherence slab. Its
phase follows the \emph{relative} motion of the pair:
\begin{equation}
  \vq_+\cdot\vr_j - \vq_-\cdot\vr_{j'} =
  \vq\cdot(\vr_j - \vr_{j'}) + \tfrac{1}{2}\Delta\vq_\perp\cdot(\vr_j + \vr_{j'}).
\end{equation}
A common axial displacement cancels exactly. A common lateral drift only
survives through the small offset $\Delta\vq_\perp \sim k\,\mathrm{NA}$. This
is the origin of the instrument's vibration immunity.

\paragraph{Reduced model.}
Eq.~\eqref{eq:Splus} can be evaluated directly from the LAMMPS positions,
without simulating the hologram. The full chain
(Eqs.~\eqref{eq:Ui}--\eqref{eq:camera}) is only needed to study pixel,
exposure, noise and aliasing effects.

% =====================================================================
\section{Fluctuation spectrum}
\label{sec:spectrum}

\subsection{Observable}

The EAFIT processing \cite[Eqs.~2.33--2.36 and \S3.1.6]{lopez2026} takes the
temporal FFT of the \emph{complex} reconstructed field $S_+$ at each pixel,
which by the Wiener--Khinchin theorem is the Fourier transform of the field
autocorrelation
\begin{equation}
  G_1(\Delta u,\Delta v,\tau) = \avg{S_+(\Delta u,\Delta v,t)\,S_+^*(\Delta u,\Delta v,t+\tau)}_t .
\end{equation}
The per-pixel power spectral density is summed over a region $\Omega$ of the
sideband and normalized by its mean over frequency:
\begin{equation}
  \Phi_\Omega(f_t) = \frac{\displaystyle\sum_{(\Delta u,\Delta v)\in\Omega}
          \left|\FT_t\{S_+(\Delta u,\Delta v,t)\}(f_t)\right|^2}
         {\displaystyle\avg{\sum_{(\Delta u,\Delta v)\in\Omega}
          \left|\FT_t\{S_+\}\right|^2}_{f_t}},
  \qquad f_t \in \left[\tfrac{1}{T},\ \tfrac{f_s}{2}\right],
  \label{eq:psd}
\end{equation}
where $\FT_t$ is the one-dimensional Fourier transform over time and $T$ is
the total acquisition time. Goodwin et al.\ and Nolte et al.\
\cite{goodwin2025,nolte2011} instead take the spectrum of the reconstructed
intensity $|S_+|^2$. The two choices give different knee frequencies, as shown
next.

\subsection{Brownian globules}

For a globule with diffusion coefficient $D_j$, the phase increment over a lag
$\tau$ is Gaussian with $\avg{(\vq\cdot\Delta\vr_j)^2} = 2q^2D_j\tau$. This is
the single-scattering limit of Ref.~\cite[Eq.~2.37]{lopez2026}. For independent
globules, the phase of a pair term in Eq.~\eqref{eq:Splus} depends on the
relative displacement, and its field correlation is
\begin{equation}
  g_1^{(jj')}(\tau) = \exp\!\left[-q^2(D_j + D_{j'})|\tau|\right].
\end{equation}
Terms from different pairs are uncorrelated, so the field autocorrelation of a
pixel is the weighted sum
\begin{equation}
  G_1(\tau) = \sum_{j,j'} |w_{jj'}|^2\,\ee^{-q^2(D_j + D_{j'})|\tau|},
  \qquad
  w_{jj'} = a_s(a_j)\,a_s^*(a_{j'})\,T_jT_{j'}\,C(\cdot)\,|\gamma(\Delta L_{jj'})|.
\end{equation}
\paragraph{Field spectrum (EAFIT observable).}
Each exponential transforms into a Lorentzian,
$\FT_t\{\ee^{-\Gamma|\tau|}\} = 2\Gamma/[\Gamma^2 + (2\pi f_t)^2]$, with its knee
at $2\pi f_t = \Gamma$. For equal globules, $\Gamma = 2q^2D$ and
\begin{equation}
  \boxed{\; f_{t,\mathrm{knee}} = \frac{q^2D}{\pi} \qquad \text{(spectrum of the field } S_+\text{)} \;}
  \label{eq:knee}
\end{equation}
This result does not need any assumption about the statistics of $S_+$.

\paragraph{Intensity spectrum (Goodwin/Nolte observable).}
If the spectrum is taken of $|S_+|^2$ instead, the result depends on how many
pairs a pixel sums. When it sums many independent pairs (the field of view
holds many resolution cells), $S_+$ is a circular Gaussian variable, the
Siegert relation gives $\avg{I(t)I(t+\tau)}/\avg{I}^2 - 1 = |G_1(\tau)/G_1(0)|^2$,
and the knee doubles to $2q^2D/\pi$. When a single product $U_aU_b^*$
dominates, the intensity correlation is $2x + x^2$ with $x = \ee^{-2q^2D\tau}$,
a mixture of Lorentzians at $q^2D/\pi$ and $2q^2D/\pi$, with an effective
knee $\approx 1.24\,q^2D/\pi$. This was checked numerically in the
sample-simulation work.

For a polydisperse emulsion, the spectrum is a superposition of Lorentzians
weighted by $|w_{jj'}|^2$, which is dominated by the large globules.

\subsection{Link to the measured knees}

L\'opez Zapata measured provolone on the EAFIT setup over 24 hours
\cite[\S4.2.2]{lopez2026}. The spectra show a knee at
$f_{t,\mathrm{knee}} \approx 5\times10^{-3}$--$10^{-2}\,$Hz that moves to lower
frequencies as the cheese dries, plus excess power below $1\,$mHz. With
Eq.~\eqref{eq:knee} and $q \approx 20.5\um^{-1}$, the knee corresponds to
\begin{equation}
  D = \frac{\pi f_{t,\mathrm{knee}}}{q^2} \approx 3.7\times10^{-5}\text{ to }
  7.5\times10^{-5}\um^2/\mathrm{s}.
\end{equation}
Ref.~\cite{lopez2026} estimates $D \sim 1.8\times10^{-3}\um^2/\mathrm{s}$
from $v = \lambda f_\mathrm{knee}$ and $D \sim v^2/(2f_\mathrm{knee})$. That
relation does not include the scattering vector, which is why the two values
differ.

Goodwin et al.\ report knees at $3\,$mHz and $30\,$mHz \cite{goodwin2025}.
If those spectra are of $|S_+|^2$ with Gaussian statistics, they correspond to
$D = \pi f/(2q^2) \approx 1.1\times10^{-5}$ and $1.1\times10^{-4}\um^2/\mathrm{s}$
(with $q \approx 20.4\um^{-1}$ at $839.6\nm$ and $30^\circ$).

For the simulations, these values set the target: the globule diffusion
coefficients in the LAMMPS model should fall around
$10^{-5}$--$10^{-4}\um^2/\mathrm{s}$, and the run should follow the EAFIT
protocol, with a frame every 5 s for about 7000 s.

% =====================================================================
\section{Open points}
\label{sec:open}

\begin{enumerate}[label=(O\arabic*)]
  \item Illumination angle $\psi$ of the EAFIT setup. Fig.~3.9 of
        Ref.~\cite{lopez2026} draws it oblique, but the text says illumination
        and collection share an axis. The angle sets the direction of $\vq$
        and the tilt of the coherence gate, and whether the specular
        reflection from the plate bottom reaches the camera.
  \item Diameters of the objective, of the field stop at the IP and of the
        aperture before FL. They set the NA (resolution cell, globules per
        cell, $\Delta\vq_\perp$), the field of view $L$ and the beam diameter
        $D$ at the biprism.
  \item Distance $d$ from the biprism to the camera, which sets the shear
        $\eta_s$ and decides whether the half-plane approximation holds.
  \item Exposure time and ROI used for the provolone acquisitions, and the
        plate (material and bottom thickness).
  \item Structure-factor correction to $\mu_s$ at $\phi\approx0.3$.
  \item Multiple scattering (deferred), which shifts the knees upward and
        makes the spectrum non-Lorentzian.
\end{enumerate}

% =====================================================================
\begin{thebibliography}{9}

\bibitem{goodwin2025}
G.~Goodwin, D.~Lim, S.~Mahajan, and D.~D.~Nolte,
``Fresnel biprism common-path low-coherence digital holography for dynamic
light scattering spectroscopy of biological materials,''
\emph{Biomed.\ Opt.\ Express} \textbf{16}(2), 806--819 (2025).
doi:10.1364/BOE.543308

\bibitem{lopez2026}
A.~L\'opez Zapata,
``Experimental and instrumental framework for dual channel angular
multiplexing polarization holography (DCAMPH) and Doppler spectroscopy,''
undergraduate thesis, EAFIT University, Medell\'in (2026).

\bibitem{nolte2011}
D.~D.~Nolte, R.~An, J.~Turek, and K.~Jeong,
``Holographic tissue dynamics spectroscopy,''
\emph{J.\ Biomed.\ Opt.} \textbf{16}(8), 087004 (2011).

% TODO: verify volume/pages of this reference
\bibitem{stocker2016}
S.~Stocker, F.~Foschum, P.~Krauter, F.~Bergmann, A.~Hohmann, C.~Scalfi Happ,
and A.~Kienle,
``Broadband optical properties of milk,''
\emph{Appl.\ Spectrosc.} \textbf{71}(5), 951--962 (2017).

\bibitem{goodman}
J.~W.~Goodman, \emph{Introduction to Fourier Optics}, 3rd ed.\
(Roberts \& Company, 2005).

\end{thebibliography}

\end{document}
